\chapter{调度方法与示例验证}

\section{滚动调度流程}

滚动调度的基本思想是在固定时间窗口内反复求解局部优化问题，并将最新观测反馈给下一轮模型。该流程适合处理天气变化、负荷波动和设备状态漂移。本文示例将流程拆分为预测、约束更新、优化求解、策略下发和结果复核五个阶段。

\begin{enumerate}
  \item 读取当前时刻的设备、环境和负荷数据；
  \item 预测未来窗口内的需求变化；
  \item 根据安全边界和运行规则更新约束；
  \item 求解调度策略并输出候选方案；
  \item 比较执行结果与预测结果，更新下一轮输入。
\end{enumerate}

\section{结果分析}

为了展示长表格和多列数据的排版效果，表 \ref{tab:hit-scheduling-result} 给出一组公开演示数据。数据仅用于模板验证，不代表真实工程系统。

\begin{table}[htbp]
  \centering
  \caption{滚动调度结果示例}
  \label{tab:hit-scheduling-result}
  \begin{tabular}{cccccc}
    \toprule
    时段 & 预测负荷/kW & 实际负荷/kW & 供能量/kWh & 峰值惩罚 & 风险评分\\
    \midrule
    1 & 128.4 & 130.1 & 126.0 & 0.12 & 0.08\\
    2 & 135.6 & 134.8 & 132.7 & 0.15 & 0.09\\
    3 & 142.2 & 143.0 & 139.1 & 0.18 & 0.11\\
    4 & 138.5 & 137.9 & 135.6 & 0.14 & 0.10\\
    \bottomrule
  \end{tabular}
\end{table}

结果表明，示例调度流程能够在保持负荷跟踪的同时控制峰值惩罚。对于真实博士论文，作者还应补充数据来源、模型参数、实验平台、消融实验和统计显著性分析。

\section{模板验证说明}

本章重点并非提出新的能源系统算法，而是验证 HIT 博士论文模板的可维护性。当前实现为 `bensz-thesis` 新增独立的 `thesis-hit-doctor` profile 和 style，项目层只维护正文、元数据和示例材料。这样的分层能避免修改哈尔滨工业大学模板时影响南京大学、中国科学院大学、中山大学等既有模板。
